\section{Mass at an isolated holomorphic zero}
\label{sec:holomorphic-mass}

The next estimate turns the order of a holomorphic zero into an integral
bound for its Jacobian minors. Its application will use high powers of the
inverse planar map. In complex dimension $d$, the mass bound grows as the
$d$th power of their exponent, matching the radial rescaling that will produce boundary
feasibility probabilities in Section~\ref{sec:feasibility}.
The estimate is a special case of the sublevel-mass and comparison
formulas for generalized Lelong numbers; see
Demailly~\cite[Sections~3 and~5]{Demailly}. We give a direct Stokes
proof in the normalization used here.

On $\mathbb C^d$, write $z_j=x_j+iy_j$ and use the complex orientation
\[
dV=dx_1\wedge dy_1\wedge\cdots\wedge dx_d\wedge dy_d.
\]
For a smooth real function $v$, our convention is
\begin{equation}\label{eq:mass-dc-convention}
d^c=\frac{i}{4}(\bar\partial-\partial),
\qquad dd^c=\frac{i}{2}\partial\bar\partial.
\end{equation}
Thus $d^cv(\xi)=-\tfrac14dv(i\xi)$ for a real tangent vector $\xi$, and
\begin{equation}\label{eq:mass-volume-normalization}
dd^c|z|^2=\sum_{j=1}^d dx_j\wedge dy_j,
\qquad
(dd^cv)^d=d!\det(v_{\bar z_i z_j})\,dV.
\end{equation}
The second identity follows by expanding the exterior product, or by a
unitary diagonalization of the Hermitian Hessian.

\Needspace{8\baselineskip}
\begin{lemma}[Holomorphic mass]\label{lem:holomorphic-mass}
Let $d\geq1$, $N\geq d$, and let $\mathcal U\subset\mathbb C^d$ be open
and contain $0$. Suppose that $f=(f_1,\ldots,f_N):\mathcal U\to\mathbb C^N$
is holomorphic and has the following properties:
\begin{enumerate}
\item $f^{-1}(0)=\{0\}$;
\item for an integer $k\geq1$,
\[
f(z)=H(z)+O(|z|^{k+1})\quad\text{near }0,
\]
where each component of $H$ is a homogeneous polynomial of degree $k$
and $H(z)\ne0$ whenever $z\ne0$;
\item for every $0<s<1$, the set
$\{z\in\mathcal U:|f(z)|^2\leq s\}$ is compact in $\mathcal U$.
\end{enumerate}
Then, with $\tau=|f|^2$,
\begin{equation}\label{eq:holomorphic-minor-mass}
\int_{\{\tau<1\}}
 \sum_{\substack{I\subset\{1,\ldots,N\}\\|I|=d}}
 \left|\det\left(\frac{\partial f_i}{\partial z_j}\right)_{
                       i\in I,\ 1\leq j\leq d}\right|^2\,dV
\geq \frac{(\pi k)^d}{d!}.
\end{equation}
All norms are Euclidean. Rows in each determinant are in increasing index
order.
\end{lemma}

\begin{proof}
We first bound the integral of $(dd^c\tau)^d$. The conversion to
Jacobian minors will then follow from Cauchy--Binet.
Both $\tau$ and $\log\tau$ away from $0$ have positive semidefinite
complex Hessians. Indeed, if $Z$ is the complex derivative matrix of $f$,
the Hessian of $\tau$ is $Z^*Z$. For a complex direction $\xi$, the
Hessian of $\log\tau$ evaluates to
\[
\frac{|f|^2|df(\xi)|^2
       -|\langle df(\xi),f\rangle|^2}{|f|^4}\geq0
\]
by Cauchy--Schwarz. Their top exterior powers are therefore nonnegative.

\paragraph{Comparison with a small sphere.}
Fix a regular value $s\in(0,1)$ of $\tau$ and set
$E_s=\{z\in\mathcal U:\tau(z)<s\}$. By compactness of the sublevel and
the regular-value theorem, $\overline{E_s}$ is a compact smooth manifold
with boundary $\{\tau=s\}$. Stokes' theorem gives
\[
\int_{E_s}(dd^c\tau)^d
 =\int_{\partial E_s}d^c\tau\wedge(dd^c\tau)^{d-1}.
\]
Put $\gamma=d^c\log\tau$ on $\mathcal U\setminus\{0\}$. The identities
\[
d^c\log\tau=\tau^{-1}d^c\tau,
\qquad
dd^c\log\tau=\tau^{-1}dd^c\tau
                 -\tau^{-2}d\tau\wedge d^c\tau
\]
and the vanishing of the pullback of $d\tau$ to $\partial E_s$ imply
\begin{equation}\label{eq:mass-level-flux}
s^{-d}\int_{E_s}(dd^c\tau)^d
 =\int_{\partial E_s}\gamma\wedge(d\gamma)^{d-1}.
\end{equation}

For sufficiently small $\varepsilon>0$, the closed Euclidean ball
$\overline{B_\varepsilon}$ is contained in $E_s$. Orient
$S_\varepsilon=\partial B_\varepsilon$ outward from this ball. Applying
Stokes on $E_s\setminus\overline{B_\varepsilon}$ gives
\begin{equation}\label{eq:mass-punctured-flux}
s^{-d}\int_{E_s}(dd^c\tau)^d
 -\int_{S_\varepsilon}\gamma\wedge(d\gamma)^{d-1}
 =\int_{E_s\setminus\overline{B_\varepsilon}}
                     (dd^c\log\tau)^d\geq0.
\end{equation}
Thus it remains to determine the flux through a shrinking sphere.

\paragraph{The homogeneous flux.}
Let $T(z)=|H(z)|^2$ and $S=\{z:|z|=1\}$. The holomorphic Taylor
expansion implies
\[
\varepsilon^{-2k}\tau(\varepsilon z)\longrightarrow T(z)
\quad\text{in }C^2\text{ on a fixed closed annulus containing }S.
\]
For example, termwise differentiation of the convergent Taylor series
gives this convergence for the rescaled map and then for its squared
norm. Since $T$ is strictly positive on that annulus, the logarithms
converge in $C^2$ as well. Pulling back to $S$ by dilation, the additive
constant $2k\log\varepsilon$ disappears under $d^c$. Consequently
\begin{equation}\label{eq:mass-homogeneous-limit}
\lim_{\varepsilon\downarrow0}
 \int_{S_\varepsilon}\gamma\wedge(d\gamma)^{d-1}
 =\int_S\alpha\wedge(d\alpha)^{d-1},
\qquad \alpha=d^c\log T\big|_S.
\end{equation}

We compare this flux with that of the radial function $|z|^{2k}$.
Homogeneity gives
\[
T(rz)=r^{2k}T(z)\quad(r>0),\qquad T(e^{i\theta}z)=T(z).
\]
Thus $\log T-k\log|z|^2$ is unchanged by dilation and by common phase
rotation. These invariances will make its contribution to the flux
vanish. On $S$, put
\[
\alpha_0=d^c\log|z|^2\big|_S,
\qquad
\beta=d^c\bigl(\log T-k\log|z|^2\bigr)\big|_S,
\qquad
\alpha_t=k\alpha_0+t\beta\quad(0\leq t\leq1).
\]
Thus $\alpha_1=\alpha$.
For the nowhere-zero tangent vector field $V(z)=iz$ on $S$, the dilation
identity and \eqref{eq:mass-dc-convention} give
$\beta(V)=0$ and $\alpha_0(V)=1/2$. The phase-rotation identity makes every
$\alpha_t$ invariant under the flow of $V$. Cartan's identity therefore
gives
\[
\iota_Vd\alpha_t
 =\mathcal L_V\alpha_t-d\bigl(\alpha_t(V)\bigr)=0,
\]
where $\iota_V$ denotes contraction and $\mathcal L_V$ the Lie derivative.
It follows that the top-degree form
$\beta\wedge(d\alpha_t)^{d-1}$ on the $(2d-1)$-dimensional sphere
vanishes: its contraction with the nonzero vector $V$ is zero.
Differentiation and integration by parts on the closed sphere now yield
\[
\frac{d}{dt}\int_S\alpha_t\wedge(d\alpha_t)^{d-1}
 =d\int_S\beta\wedge(d\alpha_t)^{d-1}=0.
\]
For $d=1$ this is simply the derivative of $\int_S\alpha_t$; the same
argument applies.

On the unit sphere, $\alpha_0$ and $d\alpha_0$ are the restrictions of
$d^c|z|^2$ and $dd^c|z|^2$, respectively. Hence a final use of Stokes and
\eqref{eq:mass-volume-normalization} gives
\begin{align*}
\int_S\alpha\wedge(d\alpha)^{d-1}
 &=k^d\int_S\alpha_0\wedge(d\alpha_0)^{d-1}\\
 &=k^d\int_{\{|z|<1\}}(dd^c|z|^2)^d\\
 &=k^d d!\,\operatorname{vol}_{2d}(\{|z|<1\})
  =(\pi k)^d.
\end{align*}

\paragraph{The minor bound.}
Equations~\eqref{eq:mass-punctured-flux} and
\eqref{eq:mass-homogeneous-limit} imply
\[
\int_{E_s}(dd^c\tau)^d\geq s^d(\pi k)^d
\]
for every regular $s\in(0,1)$. Sard's theorem supplies an increasing
sequence of such values tending to $1$. Since the integrand is
nonnegative, monotone convergence gives
\[
\int_{\{\tau<1\}}(dd^c\tau)^d\geq(\pi k)^d.
\]
Finally, the complex Hessian $Z^*Z$ and Cauchy--Binet give the identity
\[
(dd^c\tau)^d
 =d!\sum_{|I|=d}|\det Z_I|^2\,dV.
\]
Division by $d!$ proves \eqref{eq:holomorphic-minor-mass}.
\end{proof}
